\documentclass[12pt,a4paper]{article}

% --- pacotes ---
\usepackage[utf8]{inputenc}
\usepackage[portuguese]{babel}
\usepackage{geometry}
\usepackage{graphicx}
\usepackage{listings}
\usepackage{xcolor}
\usepackage{hyperref}
\usepackage{titlesec}
\usepackage{fancyhdr}
\usepackage{enumitem}
\usepackage{booktabs}
\usepackage{float}

\geometry{margin=2.5cm}

% --- estilo do codigo C ---
\lstdefinestyle{cstyle}{
  language=C,
  basicstyle=\ttfamily\small,
  keywordstyle=\color{blue}\bfseries,
  commentstyle=\color{gray}\itshape,
  stringstyle=\color{red},
  numbers=left,
  numberstyle=\tiny\color{gray},
  breaklines=true,
  frame=single,
  backgroundcolor=\color{gray!10},
  tabsize=2
}

% --- cabecalho e rodape ---
\pagestyle{fancy}
\fancyhf{}
\rhead{Programação I — 2025/2026}
\lhead{Sistema de Gestão de Workshop}
\rfoot{\thepage}

% --- titulo ---
\title{
  \vspace{-1cm}
  \Large \textbf{Escola Superior de Tecnologia e Gestão — IPVC} \\[0.5cm]
  \large Programação I — Turno PL1 \\[0.3cm]
  \rule{\linewidth}{0.5pt} \\[0.5cm]
  \LARGE \textbf{Sistema de Gestão de Participações num Workshop} \\[0.3cm]
  \rule{\linewidth}{0.5pt} \\[0.5cm]
  \large Relatório Final
}
\author{
  Artur Salgado \\ EI33385
}
\date{Junho de 2026}

% =============================================================
\begin{document}
% =============================================================

\maketitle
\thispagestyle{empty}
\newpage

% --- indice ---
\tableofcontents
\newpage

% =============================================================
\section{Introdução}
% =============================================================

[REESCREVE ESTA SECÇÃO COM AS TUAS PRÓPRIAS PALAVRAS]

O presente relatório descreve o desenvolvimento de um sistema em linguagem C para gerir as participações num workshop organizado pela ESTG. O workshop contempla um total de 16 apresentações e permite dois tipos de inscrição: inscrição aberta, que dá acesso a todo o evento, e inscrição específica, em que o participante escolhe as apresentações a que quer assistir, pagando a soma dos respetivos preços.

O sistema foi desenvolvido de acordo com os requisitos definidos no enunciado e aplica os conceitos lecionados na unidade curricular de Programação I, nomeadamente: listas ligadas, algoritmos de ordenação e pesquisa, e persistência de dados através de ficheiros binários e de texto.

% =============================================================
\section{Estruturas de Dados}
% =============================================================

\subsection{Justificação do uso de Listas Ligadas}

[REESCREVE ESTA SECÇÃO COM AS TUAS PRÓPRIAS PALAVRAS]

Optou-se pelo uso de listas ligadas para gerir os dados em memória durante a execução do programa. Esta escolha justifica-se pela flexibilidade que as listas ligadas oferecem: ao contrário dos arrays, não é necessário definir um tamanho fixo à partida, o que permite inserir e remover elementos de forma eficiente sem desperdiçar memória.

Para cada entidade do sistema foi criado um nó que encapsula os dados e um ponteiro para o elemento seguinte da lista.

\subsection{Struct Apresentacao}

\begin{lstlisting}[style=cstyle, caption={Struct Apresentacao em apresentacao.h}]
typedef struct {
  int id;
  char titulo[100];
  char orador[100];
  char curriculo[500];
  char palavras_chave[5][50];
  int duracao;       // em minutos
  float preco;
  char data[30];     // data/hora do horario
  int num_inscritos;
} Apresentacao;

typedef struct NoApresentacao {
  Apresentacao Apresentacao;
  struct NoApresentacao *proximo;
} NoApresentacao;
\end{lstlisting}

\subsection{Struct Participante}

\begin{lstlisting}[style=cstyle, caption={Struct Participante em participante.h}]
typedef struct {
  int id;
  char nome[100];
  char instituicao[100];
  int tipo_inscricao;      // 0 = aberta, 1 = especifica
  int apresentacoes[16];   // ids das apresentacoes escolhidas
  int num_apresentacoes;
  float total_pagar;       // calculado automaticamente
} Participante;
\end{lstlisting}

\subsection{Struct Utilizador}

\begin{lstlisting}[style=cstyle, caption={Struct Utilizador em utilizador.h}]
typedef struct {
  char username[50];
  char password[50];
  int perfil; // 0 = admin, 1 = participante
} Utilizador;
\end{lstlisting}

\subsection{Struct Workshop}

\begin{lstlisting}[style=cstyle, caption={Struct Workshop em workshop.h}]
typedef struct {
  char data_inicio[20];
  char data_fim[20];
  float preco_inscricao_aberta;
} Workshop;
\end{lstlisting}

% =============================================================
\section{Organização do Código}
% =============================================================

\subsection{Estrutura de ficheiros}

O projeto está organizado da seguinte forma:

\begin{verbatim}
ProjPL1_EI33385/
├── include/           # Headers (.h) - structs e prototipos
│   ├── apresentacao.h
│   ├── participante.h
│   ├── utilizador.h
│   ├── workshop.h
│   └── ficheiros.h
├── src/               # Codigo fonte (.c)
│   ├── main.c
│   ├── apresentacao.c
│   ├── participante.c
│   ├── utilizador.c
│   ├── workshop.c
│   └── ficheiros.c
├── relatorios/        # Ficheiros de texto (relatorios)
└── relatorio_final/   # Relatorio PDF
\end{verbatim}

\subsection{Separação por módulos}

O código foi dividido em módulos, cada um com o seu ficheiro \texttt{.h} (declarações) e \texttt{.c} (implementação):

\begin{table}[H]
\centering
\begin{tabular}{ll}
\toprule
\textbf{Módulo} & \textbf{Responsabilidade} \\
\midrule
\texttt{apresentacao} & Gestão das apresentações \\
\texttt{participante} & Gestão dos participantes e inscrições \\
\texttt{utilizador} & Autenticação e registo de utilizadores \\
\texttt{workshop} & Configuração do workshop \\
\texttt{ficheiros} & Persistência de dados (leitura/escrita) \\
\texttt{main} & Menu principal e fluxo do programa \\
\bottomrule
\end{tabular}
\caption{Módulos do sistema}
\end{table}

% =============================================================
\section{Funcionalidades Implementadas}
% =============================================================

\subsection{Autenticação}

[REESCREVE COM AS TUAS PALAVRAS]

No arranque, o programa pede o nome de utilizador e a password. O utilizador é procurado na lista de utilizadores carregada do ficheiro binário. Se não existir, é perguntado se pretende registar-se. Por defeito, existe sempre um utilizador \textit{admin} com a password \textit{admin} e perfil de administrador.

\subsection{Menu Administrador}

O administrador tem acesso às seguintes funcionalidades:

\subsubsection{Gestão de Participantes}
\begin{itemize}
  \item Inserir participante com tipo de inscrição (aberta ou específica)
  \item Listar todos os participantes
  \item Remover participante por ID
  \item Ver informação de um participante específico
  \item Listar participantes com inscrição aberta
  \item Listar participantes de uma apresentação
  \item Somar o total de todas as inscrições
\end{itemize}

\subsubsection{Gestão de Apresentações}
\begin{itemize}
  \item Inserir, remover e alterar apresentações
  \item Listar todas as apresentações
  \item Definir data/hora (escalonar) de uma apresentação
  \item Ordenar apresentações por título (bubble sort)
  \item Listar apresentações por data
  \item Saber qual a apresentação com mais inscrições
  \item Imprimir o programa do workshop para ficheiro de texto
\end{itemize}

\subsection{Menu Participante}

\begin{itemize}
  \item Consultar a sua inscrição
  \item Ver a agenda do workshop (por data)
  \item Ver todas as apresentações
  \item Consultar informação sobre uma apresentação específica
\end{itemize}

% =============================================================
\section{Algoritmos Relevantes}
% =============================================================

\subsection{Algoritmo de Ordenação — Bubble Sort}

[REESCREVE COM AS TUAS PALAVRAS]

Para ordenar as apresentações por título, foi implementado o algoritmo Bubble Sort. O algoritmo percorre a lista repetidamente, comparando os títulos de nós adjacentes com a função \texttt{strcmp}, e trocando os dados dos nós quando necessário. O processo repete-se até não haver mais trocas.

A principal vantagem desta abordagem é a sua simplicidade. Para uma lista de no máximo 16 apresentações, a eficiência não é um fator crítico.

\begin{lstlisting}[style=cstyle, caption={Bubble Sort por título}]
do {
  trocou = 0;
  atual = lista;
  while (atual->proximo != NULL) {
    if (strcmp(atual->Apresentacao.titulo,
               atual->proximo->Apresentacao.titulo) > 0) {
      // trocar os dados dos dois nos
      Apresentacao temp = atual->Apresentacao;
      atual->Apresentacao = atual->proximo->Apresentacao;
      atual->proximo->Apresentacao = temp;
      trocou = 1;
    }
    atual = atual->proximo;
  }
} while (trocou);
\end{lstlisting}

O mesmo algoritmo é aplicado para ordenar por data. Como o formato usado é \texttt{YYYY-MM-DD HH:MM}, a comparação direta de strings com \texttt{strcmp} produz a ordenação cronológica correta.

\subsection{Cálculo do Total a Pagar}

[REESCREVE COM AS TUAS PALAVRAS]

Quando o administrador inscreve um participante com inscrição específica, o sistema percorre as apresentações escolhidas e soma os respetivos preços:

\begin{lstlisting}[style=cstyle, caption={Cálculo do total na inscrição específica}]
NoApresentacao *atual = listaApresentacoes;
while (atual != NULL && atual->Apresentacao.id != id) {
  atual = atual->proximo;
}
if (atual != NULL) {
  p.apresentacoes[p.num_apresentacoes++] = id;
  p.total_pagar += atual->Apresentacao.preco;
  atual->Apresentacao.num_inscritos++;
}
\end{lstlisting}

\subsection{Persistência com Ficheiros Binários}

[REESCREVE COM AS TUAS PALAVRAS]

Para guardar e carregar os dados entre sessões, foram utilizadas as funções \texttt{fwrite} e \texttt{fread} da biblioteca padrão do C. Cada struct é escrita diretamente em binário, byte a byte. Na leitura, os dados são reconstruídos na lista ligada.

\begin{lstlisting}[style=cstyle, caption={Guardar participantes em ficheiro binário}]
while (atual != NULL) {
  fwrite(&atual->participante, sizeof(Participante), 1, fp);
  atual = atual->proximo;
}
\end{lstlisting}

\begin{lstlisting}[style=cstyle, caption={Carregar participantes do ficheiro}]
while (fread(&p, sizeof(Participante), 1, fp) == 1) {
  NoParticipante *novo = malloc(sizeof(NoParticipante));
  novo->participante = p;
  novo->proximo = *lista;
  *lista = novo;
}
\end{lstlisting}

% =============================================================
\section{Problemas Encontrados e Soluções}
% =============================================================

[REESCREVE TODA ESTA SECÇÃO COM AS TUAS PALAVRAS E EXPERIÊNCIAS REAIS]

\subsection{Buffer Overflow na String de Data}

Durante o desenvolvimento, o programa apresentava um crash (trace trap) ao tentar inserir uma apresentação. Após análise, verificou-se que o campo \texttt{data[20]} era demasiado pequeno para a string \texttt{"sem horario definido"} (21 caracteres com o terminador). A solução foi aumentar o campo para \texttt{data[30]}.

\subsection{Problema com Leitura de Strings com Espaços}

Ao usar \texttt{scanf("\%s")} para ler o nome do participante, apenas a primeira palavra era lida. Por exemplo, ao escrever "Artur Salgado", apenas "Artur" era guardado. A solução foi usar \texttt{scanf(" \%[{\textasciicircum}\textbackslash n]")} que lê até ao fim da linha, incluindo espaços.

\subsection{Buffer do Terminal entre Leituras}

Após ler um inteiro com \texttt{scanf("\%d")}, o caractere \texttt{\textbackslash n} ficava no buffer e era lido pela próxima chamada. A solução foi limpar o buffer com \texttt{while(getchar() != '\textbackslash n')} antes de ler strings longas no menu de inserção de apresentações.

% =============================================================
\section{Conclusão}
% =============================================================

[REESCREVE COM AS TUAS PALAVRAS]

O sistema desenvolvido cumpre todos os requisitos definidos no enunciado. Foram aplicados os conceitos lecionados nas aulas: listas ligadas para gestão de dados em memória, algoritmos de ordenação (bubble sort) e pesquisa, ficheiros binários para persistência permanente e ficheiros de texto para relatórios.

O projeto permitiu consolidar os conhecimentos de programação em C, nomeadamente a gestão de memória dinâmica com \texttt{malloc} e \texttt{free}, a utilização de structs e ponteiros, e a organização modular do código.

% =============================================================
\appendix
\section{Manual do Utilizador}
% =============================================================

\subsection{Compilação e Execução}

Para compilar o programa, executar o seguinte comando na pasta raiz do projeto:

\begin{verbatim}
gcc src/*.c -I include/ -o programa
\end{verbatim}

Para executar:

\begin{verbatim}
./programa
\end{verbatim}

\subsection{Login}

Ao iniciar o programa, é pedido o nome de utilizador e a password.

\begin{itemize}
  \item \textbf{Admin por defeito:} utilizador \texttt{admin}, password \texttt{admin}
  \item Se o utilizador não existir, o programa pergunta se pretende registar-se
  \item Um novo utilizador registado tem sempre perfil de \textit{participante}
\end{itemize}

\subsection{Menu Administrador}

\begin{table}[H]
\centering
\begin{tabular}{cl}
\toprule
\textbf{Opção} & \textbf{Descrição} \\
\midrule
1  & Inserir participante (com tipo de inscrição) \\
2  & Listar todos os participantes \\
3  & Remover participante por ID \\
4  & Ver informação de um participante específico \\
5  & Listar participantes com inscrição aberta \\
6  & Listar participantes de uma apresentação \\
7  & Somar total de todas as inscrições \\
8  & Inserir nova apresentação \\
9  & Listar todas as apresentações \\
10 & Remover apresentação por ID \\
11 & Alterar dados de uma apresentação \\
12 & Definir data/hora de uma apresentação \\
13 & Ordenar apresentações por título \\
14 & Listar apresentações por data \\
15 & Apresentação com mais inscrições \\
16 & Imprimir programa para ficheiro de texto \\
17 & Configurar datas e preço do workshop \\
0  & Sair (guarda todos os dados automaticamente) \\
\bottomrule
\end{tabular}
\caption{Opções do menu administrador}
\end{table}

\subsection{Menu Participante}

\begin{table}[H]
\centering
\begin{tabular}{cl}
\toprule
\textbf{Opção} & \textbf{Descrição} \\
\midrule
1 & Consultar a minha inscrição \\
2 & Ver agenda do workshop (por data) \\
3 & Ver todas as apresentações \\
4 & Consultar informação sobre uma apresentação \\
0 & Sair \\
\bottomrule
\end{tabular}
\caption{Opções do menu participante}
\end{table}

\subsection{Persistência de Dados}

Os dados são guardados automaticamente ao sair pelo menu do administrador (opção 0). São criados os seguintes ficheiros na pasta de execução do programa:

\begin{itemize}
  \item \texttt{PARTICIPANTES} — dados dos participantes (binário)
  \item \texttt{APRESENTACOES} — dados das apresentações (binário)
  \item \texttt{UTILIZADORES} — dados dos utilizadores (binário)
  \item \texttt{WORKSHOP} — configuração do workshop (binário)
  \item \texttt{relatorios/programa\_workshop.txt} — programa do workshop (texto)
\end{itemize}

\subsection{Tipo de Inscrição}

Ao inserir um participante (opção 1 do admin):
\begin{itemize}
  \item \textbf{0 — Inscrição Aberta:} acesso total ao workshop, paga o preço configurado em "Configurar Workshop"
  \item \textbf{1 — Inscrição Específica:} escolhe as apresentações pelos IDs (são listadas automaticamente), o total é a soma dos preços das apresentações escolhidas
\end{itemize}

\end{document}
